\section{Introduction}
\label{sec:intro}

% CO-AUTHOR EDIT: REWRITE OPENING PARAGRAPH [T3.1]
Bid-rigging cartels in public procurement rely on a specific mechanism to
avoid detection: \emph{cover bidding}---the deployment of firms that
submit deliberately losing bids to simulate competition and satisfy
minimum-bidder requirements \citep{marshall2012economics}. We call these
firms \emph{frequent losers} (FL): firms that never win a single tender
yet participate in an abnormally large number of procurement processes.
Using 4.5 million tender-items from S\~ao Paulo's electronic procurement
platform, we show that FL presence is associated with 4--9\% higher
prices, and that the FL screen achieves an AUC of 0.94 when validated
against Brazil's competition authority (CADE) cartel convictions---
substantially outperforming bid-level detection methods that require
granular price data. For procurement agencies in developing countries,
where bid-level data are often unavailable, the FL screen offers the
first firm-level, ex ante screening tool operable with participation
records alone.

Public procurement accounts for 12--15\% of GDP in OECD countries and a
substantially larger share in developing economies
\citep{oecd2019government}, and bid rigging inflates costs by an
estimated 10--50\% \citep{connor2007price}. Existing detection
approaches rely either on \emph{reactive} methods (leniency programs,
whistleblowers) or on \emph{proactive} statistical screens that
typically require granular bid-level data \citep{harrington2008detecting,
imhof2019detecting}. The FL screen complements these approaches with two
practical advantages: (i)~it requires only participation and outcome
data, not bid values; and (ii)~it produces firm-level flags that
implicate all tenders in which a flagged firm participates, enabling
targeted investigation.

This paper advances two hypotheses of increasing strength:

\begin{description}
  \item[H1 (Screening marker):] FL presence is a useful
    \emph{screening marker} for potential collusion---it correlates
    with price anomalies and overlap with known cartel participants,
    regardless of whether every FL firm is itself a cover bidder.
  \item[H2 (Causal mechanism):] FL firms are \emph{themselves} cover
    bidders deployed by cartels to inflate bidder counts---implying a
    direct causal link between FL presence and price inflation.
\end{description}

\noindent Our primary contribution is to show that FL presence is a
robust and implementable screening marker for bid rigging (H1). We then
show that a range of independent empirical tests---instrumental
variables, co-bidding network analysis, and bid-level
diagnostics---converge in a manner most consistent with a cover-bidding
mechanism (H2). We do not claim to identify all cartels in the sample,
estimate structural overcharges, or measure aggregate collusion welfare
losses. The FL screen complements existing tender-level forensic tools
\citep{imhof2019detecting} with a \emph{firm-level} criterion that is
operable with minimal data---participation records and outcome
flags---making it particularly relevant for competition authorities in
developing countries where bid-level data are unavailable.

\subsection{Contributions}

Using 4.5 million tender-items and 40 million bids from S\~ao Paulo's
Bolsa Eletr\^onica de Compras (BEC) platform (2009--2019), we make three
contributions.

\emph{First}, we show that FL presence is a reliable empirical marker % V7-T2: cross-fit elevated
for procurement anomalies (H1). FL-present tenders exhibit 4--9\%
higher negotiated prices in OLS specifications with item, year, and
purchasing-unit fixed effects (preferred specification: 6.4\%). % V8-POLISH: 0.0636 ≈ 6.4%
Temporal cross-fitting---defining FL on odd years and estimating on
even years, and vice versa---breaks any mechanical link between
classification and outcomes and yields a conservative estimate of
3.6\%, confirming that the FL--price association is not an artifact
of within-sample classification. External validation against CADE
(Brazil's competition authority) convictions shows that 7.1\% of FL
firms co-participate with convicted cartelists---3.5 times the rate
expected by chance after controlling for participation intensity
($p < 0.001$, permutation test). The FL--price association persists
after excluding all CADE-involved markets, indicating that the screen
detects anomalies beyond known cartels.

\emph{Second}, we provide evidence consistent with the causal mechanism
(H2) using a leave-one-out instrumental variable (IV) strategy. The
instrument exploits supply-side variation in FL firms' availability: for
each purchasing unit, we measure FL activity at \emph{other} purchasing
units in the same product market and year. The first stage is strong
($F = 396$) and the IV estimate of a 21.4\% price markup exceeds the OLS % V8-POLISH: exp(0.194)-1=21.4%
estimate, consistent with measurement-error attenuation. Co-bidding
network analysis reveals that the FL--price effect is concentrated in
competitive markets (low winner concentration), where cover bidding is
most profitable, rather than in concentrated markets where a dominant
firm already controls pricing.

\emph{Third}, we provide statistical evidence of bid coordination using
\cites{bajari2003deciding} framework. FL bid residuals violate both
exchangeability (Kolmogorov--Smirnov $D = 0.15$, $p < 0.001$) and
conditional independence (pairwise product = 4.28 for FL, $t = 81.0$, % R1-FIX: updated to v6 enriched first-stage results
$p < 0.001$; bootstrap 95\% CI for the FL--non-FL difference excludes
zero).

% CO-AUTHOR EDIT: CONTRIBUTION TABLE [T3.3]
Table~\ref{tab:contributions} summarizes the FL screen's positioning
relative to existing approaches.

\begin{table}[htbp]
\centering
\small
\caption{Positioning Relative to Existing Cartel Detection Methods}
\label{tab:contributions}
\begin{tabular}{lccc}
\toprule
 & \textbf{This paper} & Bid-level screens & Structural tests \\
 & (FL screen) & (Imhof et al., 2019) & (Bajari \& Ye, 2003) \\
\midrule
Unit of detection & Firm & Tender & Tender \\
Requires bid values & No & Yes & Yes \\
Scalable to large data & Yes & Partial & No \\
Validated vs.\ CADE & Yes (AUC$=$0.94) & AUC$=$0.79 & N/A \\
Ex ante screening & Yes & No & No \\
Policy threshold & Yes & No & No \\
\bottomrule
\end{tabular}
\end{table}

The remainder of this paper is organized as follows.
Section~\ref{sec:structural} develops a framework for cover-bidder
deployment with testable predictions. Section~\ref{sec:data} describes
the data and FL definition. Section~\ref{sec:cade} provides external
validation against CADE convictions. Section~\ref{sec:empirical} details
the empirical strategy. Sections~\ref{sec:results}
and~\ref{sec:mechanisms} present results and mechanism tests.
Section~\ref{sec:robustness} reports robustness checks.
Section~\ref{sec:counterfactual} derives policy counterfactuals. % CO-AUTHOR EDIT: SECTION REORDER [T3.4]
Section~\ref{sec:conclusion} concludes.
